\section{Proof of the Universal Representation Theorem}
\label{app:theory-strengthening}

We prove \cref{thm:universal-interventional-calculus} in full.

\subsection{Canonical causal realization}

Fix an intervention configuration $a$.  At time $k$, let $h_k$ denote the
complete extended history of oracle observations, exogenous inputs, realized
innovations, and previous outputs.  All component spaces are standard Borel.
Causality means that, after exposing the protocol's innovation, the current
response is a measurable function of $h_k$, the current declared input, and
that innovation, without dependence on future inputs or innovations.  Taking
$x_k^a=h_k$ and letting $\Phi_k^a$ append the next declared input,
observation, innovation, and output gives a pathwise unifilar history-state
realization of \cref{eq:causal-response}.  Thus existence requires no
finite-memory assumption.

For extended histories at the same time, define $h\sim h'$ when every common
future sequence of declared inputs and innovation values produces identical
future response sequences.  This is an equivalence relation.  Let $[h]$ be
the equivalence class.  If $h\sim h'$, append the same next input and
innovation value.  Every later input--innovation sequence is then a common
continuation of the original histories, so the two extended histories remain
equivalent.  The transition $[h]\mapsto[h^+]$ is therefore well-defined, and
the response map is constant on equivalence classes.  The classes form a
sufficient pathwise predictive-state realization in the category defined in
\cref{sec:language}.

Now let $s(h)$ be the state reached by any other admissible sufficient
pathwise/unifilar realization.  If $s(h)=s(h')$, its response map and
unifilar updates produce the same future response sequence under every common
input--innovation continuation, so $h\sim h'$.  Hence
\[
\pi:s(h)\longmapsto[h]
\]
is a well-defined surjection from its reachable states onto the predictive
quotient.  A realization with two reachable states mapped to the same class
is not behaviorally minimal.  Every behaviorally minimal reachable
realization therefore makes $\pi$ bijective.  Its transition and response
maps are conjugate to the quotient maps, proving uniqueness up to state
relabeling.  This is a set-level behavioral quotient relative to the fixed
innovation representation.  A standard-Borel structure on the quotient
requires the induced equivalence relation to be smooth; the theorem does not
assume that additional descriptive-set-theoretic property.

\subsection{Incidence-algebra completeness}

Let $\mu_{\mathcal P}$ be the M\"obius function of the lower-finite poset.
Every sum below is finite.  Define
\[
\zeta_t=\sum_{s\preceq t}\mu_{\mathcal P}(s,t)F(s).
\]
For fixed $a\in\mathcal P$, exchange the finite sums:
\begin{align*}
\sum_{t\preceq a}\zeta_t
&=\sum_{t\preceq a}\sum_{s\preceq t}
\mu_{\mathcal P}(s,t)F(s)\\
&=\sum_{s\preceq a}F(s)
\sum_{t:s\preceq t\preceq a}\mu_{\mathcal P}(s,t).
\end{align*}
By the defining inverse identity for the M\"obius function, the inner sum is
one when $s=a$ and zero otherwise.  The result is $F(a)$.  Applying the same
M\"obius transform to any proposed reconstruction proves uniqueness.

For the Boolean lattice, the interval $[B,T]$ is a Boolean lattice of rank
$|T|-|B|$, so $\mu(B,T)=(-1)^{|T|-|B|}$ and the formula reduces to mixed
finite differences.  For a product of chains, the product M\"obius function
gives the corresponding graded mixed differences.

\subsection{The complete experimental gauge}

Assume $\mathbb H\cong\mathbb R^d$ and stack the supported effect vectors.
For queried configurations $\mathcal D=(a_1,\ldots,a_n)$,
\begin{equation}
y_{\mathcal D}
=(X_{\mathcal D,\mathcal K}\otimes I_d)\zeta,
\qquad
X_{\mathcal D,\mathcal K}(i,t)=\mathbf1\{t\preceq a_i\}.
\label{eq:appendix-mask-linear-system}
\end{equation}
Two collections produce the same observations exactly when their difference
lies in the kernel.  Since
\[
\rank(X_{\mathcal D,\mathcal K}\otimes I_d)
=d\rank X_{\mathcal D,\mathcal K},
\]
rank--nullity gives compatible-fiber dimension
$d(|\mathcal K|-r_{\mathcal D})$ and proves the identification criterion.

If a finite poset has a greatest element $\top$, then $t\preceq\top$ for every
$t\in\mathcal K$.  The single passive row is all ones, has rank one, and
leaves ambiguity dimension $d(|\mathcal K|-1)$.

\subsection{Sharp quotient stability}

Put $A=X\otimes I_d$ and first suppose $X$ has positive rank.  Let
$\widehat\zeta=A^\dagger y$.  Since $A^\dagger A$ projects orthogonally onto
$\ker A^\perp$,
\[
\widehat\zeta-\zeta_{\rm can}
=A^\dagger e
=A^\dagger\operatorname{Proj}_{\operatorname{range}A}e.
\]
The positive singular values of $X\otimes I_d$ are those of $X$, each
repeated $d$ times.  Therefore
\[
\norm{\widehat\zeta-\zeta_{\rm can}}
\le
\frac{\norm{\operatorname{Proj}_{\operatorname{range}A}e}}
{\sigma_{\min}^+(X)}.
\]
Equality holds when the projected error is a left singular vector associated
with $\sigma_{\min}^+(X)$.
If $X$ has rank zero, $\ker A$ is the entire coordinate space and its
canonical identifiable quotient is $\{0\}$, so there is no positive singular
direction to estimate.

\subsection{Exact fixed and adaptive intervention complexity}

An $n\times|\mathcal K|$ matrix cannot have full column rank for
$n<|\mathcal K|$, so every fixed exact design needs at least
$|\mathcal K|$ queries.  To attain the bound, query the configurations in
$\mathcal K$.  Choose a linear extension of the partial order and use it for
both rows and columns.  In the square matrix
\[
X_{a,t}=\mathbf1\{t\preceq a\},
\qquad a,t\in\mathcal K,
\]
a nonzero entry can occur only when the column precedes the row, and every
diagonal entry is one.  The matrix is lower triangular with unit diagonal and
is invertible.

For the deterministic adaptive lower bound, run a strategy with fewer than
$|\mathcal K|$ queries against the zero effect law.  Let $\mathcal D_0$ be its
realized configurations.  Choose nonzero
$c\in\ker X_{\mathcal D_0,\mathcal K}$ and nonzero $v\in\Hspace$, and set
$h_t=c_tv$.  The zero law and the law $\zeta'=h$
return identical responses at the first query.  If their transcripts agree
through one query, determinism selects the same next configuration, where the
responses again agree by the kernel construction.  Induction gives identical
complete transcripts under distinct laws.

Now consider a randomized strategy.  Although the ambient lower-finite poset
need not be finite, a query produces one of at most $2^{|\mathcal K|}$
distinct incidence row patterns on the finite support $\mathcal K$.  Fewer
than $|\mathcal K|$ queries therefore produce one of finitely many row-pattern
sequences.  Under the zero law, at least one such sequence $R_0$ occurs with
positive probability.  Its stacked matrix has fewer than $|\mathcal K|$
rows, so choose nonzero $c$ in its kernel, fix nonzero $v\in\Hspace$, and set
$h_t=c_tv$.  On the event producing $R_0$, the zero law and $h$ return
identical responses at every queried pattern.  With
the same internal random choices, induction gives the same subsequent
configurations and hence the same full transcript.  The estimator returns
the same value for two distinct laws on an event of positive probability, so
it cannot be correct with probability one for both.  Randomization cannot
beat the lower bound.

\subsection{Prediction and falsification}

When $X_{\mathcal D,\mathcal K}$ has full column rank, the fitting responses
determine a unique $\zeta$.  The structural model forces
\[
F_{\mathcal K}(b)
=\sum_{\substack{t\in\mathcal K\\t\preceq b}}\zeta_t
\qquad\forall b\in\mathcal P.
\]
A held-out response differing from this value cannot belong to the declared
support class.  This proves falsifiability and completes the theorem.
